% problem_id: S046-thm:prod-ump
% source_id: S046 (HoTT Book)
% source_file: basics.tex
% statement_locator: basics.tex lines 2372--2375
% proof_locator: basics.tex lines 2376--2390
% env: thm
% pairing: tight (verified)
% status: complete (statement + author proof verbatim + reason + locators)
%
\begin{thm}\label{thm:prod-ump}
  \index{universal!property!of cartesian product}%
  \eqref{eq:prod-ump-map} is an equivalence.
\end{thm}

\begin{proof}
  We define the quasi-inverse by sending $(g,h)$ to $\lam{x}(g(x),h(x))$.
  (Technically, we have used the induction principle for the cartesian product $(X\to A)\times (X\to B)$, to reduce to the case of a pair.
  From now on we will often apply this principle without explicit mention.)

  Now given $f:X\to A\times B$, the round-trip composite yields the function
  \begin{equation}
    \lam{x} (\proj1(f(x)),\proj2(f(x))).\label{eq:prod-ump-rt1}
  \end{equation}
  By \cref{thm:path-prod}, for any $x:X$ we have $(\proj1(f(x)),\proj2(f(x))) = f(x)$.
  Thus, by function extensionality, the function~\eqref{eq:prod-ump-rt1} is equal to $f$.

  On the other hand, given $(g,h)$, the round-trip composite yields the pair $(\lam{x} g(x),\lam{x} h(x))$.
  By the uniqueness principle for functions, this is (judgmentally) equal to $(g,h)$.
\end{proof}
